\documentclass[10pt,conference]{IEEEtran}
\usepackage{cite}
\usepackage{amsmath,amssymb,amsfonts}
\usepackage{algorithmic}
\usepackage{graphicx}
\usepackage{textcomp}
\usepackage{xcolor}
\usepackage{booktabs}
\usepackage{url}
\usepackage{hyperref}

\def\BibTeX{{\rm B\kern-.05em{\sc i\kern-.025em b}\kern-.08em
    T\kern-.1667em\lower.7ex\hbox{E}\kern-.125emX}}

\begin{document}

\title{recovery\_eval: State-Matched and Provenance-Aware Evaluation of Recovery Behavior in Language-Model Reasoning}

\author{\IEEEauthorblockN{Sham Satish Thakare}
\IEEEauthorblockA{\textit{Independent Researcher} \\
Pune, Maharashtra, India \\
Email: shamthakare3000@gmail.com}
}

\maketitle

\begin{abstract}
Evaluating whether post-training enhances a language model's ability to recover from intermediate reasoning errors remains challenging due to confounding variables in trajectory evaluation. In this paper, we introduce \texttt{recovery\_eval}, a data-centric evaluation framework that pairs verifier-defined error states with prospectively matched reference control states using frozen structural covariates and append-only exposure governance. We validate the framework across 400 genuine neural rollouts generated by two released checkpoints (\texttt{Qwen2.5-Math-1.5B} Base and Instruct) on 20 synthetic arithmetic problems inspired by GSM8K reasoning under a corrected state-matched protocol (Protocol 2.0; 256-token budget). Under state-matched prefixes, neither model exhibits spontaneous error recovery, achieving exactly 0.00\% accuracy across all 200 recovery continuations ($p_{\text{BR}} = p_{\text{IR}} = 0.0000$). Instead, models exhibit continuation inertia: 98.5\% of recovery rollouts faithfully propagate the injected error downstream. On unperturbed control prefixes, both models achieve near-perfect accuracy (Base: 97.00\% heuristic, 100.00\% structural; Instruct: 100.00\%), yielding a matched recovery contrast of $D = 0.0000$ ($SE = 0.0000$) under structural verification and $D = -0.0300$ ($SE_{\text{cluster}} = 0.0164, p = 0.0828$) under heuristic parsing. These findings demonstrate that historical reports of a significant negative post-training interaction ($D = -0.21$) were artifacts of prompt length asymmetry, 64-token truncation censoring, and heuristic verifiers, establishing the necessity of state-matched evaluation infrastructure for reasoning diagnostics.
\end{abstract}

\begin{IEEEkeywords}
Language Models, Mathematical Reasoning, Error Recovery, State-Matched Evaluation, Data-Centric AI, Benchmark Governance, Continuation Inertia.
\end{IEEEkeywords}

\section{Introduction}
Large language models (LLMs) trained on mathematical reasoning tasks exhibit substantial improvements in step-by-step problem-solving \cite{cobbe2021gsm8k, qwen25math2024}. A foundational question in machine learning evaluation is whether post-training procedures (such as instruction tuning, process supervision, or reinforcement learning) instill structural intelligence that enables models to detect and recover from early intermediate arithmetic or logical errors \cite{lightman2023process, zelikman2022star}.

However, evaluating error recovery behavior is acutely vulnerable to evaluation artifacts. Naive comparisons between error-containing trajectories and unperturbed baselines routinely conflate trajectory depth, remaining solution length, token complexity, and prompt phrasing. Furthermore, tight generation limits induce differential truncation censoring, masking true baseline capabilities. In preliminary experiments using an unbalanced 64-token protocol (Protocol 1.0), an apparent negative interaction was observed ($D \approx -0.21$), suggesting that instruction tuning degraded recovery capability.

To resolve these ambiguities, we present \texttt{recovery\_eval}, a reproducible data-centric framework that enforces prospective caliper matching on intermediate reasoning states and maintains end-to-end cryptographic provenance. Across 400 genuine neural rollouts evaluated under a length-balanced, state-matched protocol with an expanded 256-token budget (Protocol 2.0), we establish that:
\begin{enumerate}
    \item \textbf{Dissolution of the Negative Interaction}: Under state-matched prompting, the historical negative interaction dissolves ($D = -0.0300, p = 0.0828$ under heuristic parsing; $D = 0.0000$ under structural parsing).
    \item \textbf{Continuation Inertia}: Small (1.5B) language models do not spontaneously self-correct from injected intermediate errors ($0/200$ recoveries). In 98.5\% of cases, they treat corrupted premises as authoritative and faithfully propagate errors downstream.
    \item \textbf{Evaluator Brittleness}: Heuristic last-token regex extractors artificially penalize base models by 3\% due to trailing conversational explanations, whereas structural parsing demonstrates 100\% baseline accuracy for both models.
\end{enumerate}

\section{Related Work}
Self-correction and backtracking in LLMs have attracted intense recent scrutiny \cite{snell2024scaling, wang2022selfconsistency, huang2023large}. Table \ref{tab:literature} compares \texttt{recovery\_eval} with existing paradigms.

\begin{table}[htbp]
\caption{Comparison with Existing Reasoning Evaluation Frameworks}
\label{tab:literature}
\centering
\begin{tabular}{p{2.2cm}p{1.4cm}p{1.6cm}p{1.8cm}}
\toprule
\textbf{Framework} & \textbf{State Matching} & \textbf{Exposure Ledger} & \textbf{Primitive Provenance} \\
\midrule
End-to-End GSM8K \cite{cobbe2021gsm8k} & None & No & Final Answer Only \\
PRM800K \cite{lightman2023process} & Unmatched & Partial & Step Scores \\
STaR \cite{zelikman2022star} & Unmatched & No & Filtered Traces \\
\textbf{\texttt{recovery\_eval} (Ours)} & \textbf{State-Matched} & \textbf{Append-Only LEDGER} & \textbf{Full BPE Token + Raw JSONL} \\
\bottomrule
\end{tabular}
\end{table}

Unlike end-to-end benchmarks, \texttt{recovery\_eval} prospectively matches recovery states with control states prior to generation, eliminating prefix length and depth confounds.

\section{Problem Definition \& Estimand}
Let $s$ denote a partial reasoning trajectory prefix for an arithmetic problem instance $p \in \mathcal{P}$. A deterministic verifier $\mathcal{V}(s)$ audits intermediate step correctness. A \textit{recovery state} $s_R$ contains a verifier-identified arithmetic error in Step 1. A \textit{reference control state} $s_C$ represents an unperturbed valid Step 1 prefix for the same problem at an equivalent problem-solving horizon.

For a target policy $\pi$ (e.g., Instruct) and baseline policy $\pi_0$ (e.g., Base), let $V(s) \in \{0, 1\}$ denote the binary correctness indicator of completing the solution correctly from prefix $s$. We define the matched recovery contrast $D$ as:
\begin{equation}
D = \mathbb{E}[V_{\pi}(s_R) - V_{\pi_0}(s_R)] - \mathbb{E}[V_{\pi}(s_C) - V_{\pi_0}(s_C)]
\end{equation}

\section{The \texttt{recovery\_eval} Protocol}

\subsection{State Construction and Injected Perturbation}
Recovery states $s_R$ are constructed by introducing controlled single-step arithmetic perturbations ($+2$ integer shift on intermediate multiplication products) into Step 1 prefixes. Reference control states $s_C$ maintain unperturbed Step 1 calculations.

\subsection{Prospective Caliper Matching}
To guarantee statistical balance, each recovery state $s_R$ is matched to a control state $s_C$ over structural covariates (trajectory depth, remaining length, and token length) ensuring exact categorical matching on reasoning operation (chained integer multiplication). In Protocol 2.0, Control and Recovery prefixes have identical character and token lengths ($\Delta_{\text{chars}} = 0$).

\subsection{Exposure \& Provenance Governance}
All evaluation items are governed by an append-only event ledger with SHA-256 parent hash chaining to eliminate data leakage. Raw rollout records store input token IDs, generated BPE token IDs, wall-clock generation durations, and exact model weight manifest digests.

\section{Experimental Configuration}
The empirical campaign evaluates 20 synthetic multi-step arithmetic problems inspired by GSM8K word problems ($N=20$). For each problem, 1 matched recovery state and 1 matched control state are evaluated across 2 model checkpoints using 5 stochastic generation seeds ($S=5$), yielding exactly 400 rollouts ($20 \times 2 \times 2 \times 5 = 400$).

\begin{table}[htbp]
\caption{Evaluation Design and Provenance Statistics}
\label{tab:provenance}
\centering
\begin{tabular}{lp{5.2cm}}
\toprule
\textbf{Attribute} & \textbf{Specification / Value} \\
\midrule
Base Model & \texttt{Qwen/Qwen2.5-Math-1.5B} (Rev \texttt{4a83ca6e}) \\
Instruct Model & \texttt{Qwen/Qwen2.5-Math-1.5B-Instruct} (Rev \texttt{aafeb0fc}) \\
Hardware Platform & NVIDIA GeForce RTX 4090 (24GB VRAM, CUDA 12.4) \\
Execution Precision & Native \texttt{float16} PyTorch 2.4.0 \\
Generation Config & $T=0.7$, Top-$p=0.9$, $\text{max\_new\_tokens}=256$ \\
Total Rollouts & 400 (200 Base, 200 Instruct) \\
Total Generated Tokens & 11,917 tokens (Mean: 29.79 tok/rollout) \\
Inference Duration & 265.40 seconds ($45.05\text{ tokens/sec}$) \\
Truncation Rate & \textbf{0.00\% (0 / 400 rollouts truncated)} \\
EOS Termination Rate & \textbf{100.00\% (400 / 400 rollouts hit EOS)} \\
Raw Dataset SHA-256 & \texttt{b8792f8df375fccdbd800d440cc884ef...} \\
\bottomrule
\end{tabular}
\end{table}

\section{Empirical Results}

\subsection{Cluster-Robust Factorial Outcomes}
Table \ref{tab:outcomes} summarizes completion accuracies across all four experimental cells under both heuristic and structural verification regimes.

\begin{table}[htbp]
\caption{Protocol 2.0 Factorial Outcomes ($N=400$ Rollouts, $C=20$ Clusters)}
\label{tab:outcomes}
\centering
\begin{tabular}{lccc}
\toprule
\textbf{State Condition} & \textbf{Base ($\pi_0$)} & \textbf{Instruct ($\pi$)} & \textbf{Difference ($\Delta$)} \\
\midrule
Recovery States ($s_R$) & 0.0000 & 0.0000 & $+0.0000$ \\
Control States (Heuristic) & 0.9700 & 1.0000 & $+0.0300$ \\
Control States (Structural) & 1.0000 & 1.0000 & $+0.0000$ \\
\midrule
\textbf{Matched Contrast (Heuristic)} & -- & -- & \textbf{$-0.0300$} \\
\textbf{Cluster SE / $p$-value} & -- & -- & $SE = 0.0164, p = 0.0828$ \\
\textbf{95\% Confidence Interval} & -- & -- & $[-0.0643, +0.0043]$ \\
\midrule
\textbf{Matched Contrast (Structural)} & -- & -- & \textbf{$+0.0000$} \\
\textbf{Cluster SE / $p$-value} & -- & -- & $SE = 0.0000, p = 1.0000$ \\
\bottomrule
\end{tabular}
\end{table}

\subsection{Dissolution of Negative Interaction}
Under Protocol 2.0, the historical negative interaction dissolves. Under the heuristic parser, $D = -0.0300$ ($SE_{\text{cluster}} = 0.0164, t = -1.831, p = 0.0828$). Because the 95\% confidence interval ($[-0.0643, +0.0043]$) crosses zero, the effect is statistically indistinguishable from zero at $\alpha = 0.05$. Under structural arithmetic parsing, $D = 0.0000$ exactly.

\subsection{Independent Verifier Falsification}
The 3\% deficit in Base Control (97\% vs. 100\%) was audited trace-by-trace. In all three disputed rollouts, the Base model calculated the correct revenue ($180, 176, 360$) but appended a post-solution explanation concluding in `"- In step 3..."`. The heuristic extractor's last-token fallback rule parsed the step index `3` instead of the revenue answer. Under structural equivalence, Base Control accuracy is 100/100 (1.0000).

\subsection{Behavioral Trace Analysis: Continuation Inertia}
Across all 200 recovery rollouts:
\begin{itemize}
    \item \textbf{Direct Error Propagation (98.5\%, 197/200)}: Models adopted the corrupted Step 1 quantity $E$ and computed $4 \times E$.
    \item \textbf{Attention-Slip Divergence (1.5\%, 3/200)}: In `gsm8k_fresh_004`, the model bound the unit price token into the hourly rate.
    \item \textbf{Zero Backtracking (0.0\%, 0/200)}: Zero rollouts contained error-detection words (`wait`, `actually`, `mistake`).
\end{itemize}

\subsection{Uncertainty Bounds for Zero Recovery}
While the observed recovery rate is $0/200$ ($0.0\%$), naive binomial models ($95\%$ upper bound $1.49\%$) ignore intra-cluster correlation. Accounting for $C=20$ problem clusters ($N_{\text{eff}} \approx 20$), the exact Clopper-Pearson 95\% upper bound is **13.91\%**.

\section{Limitations}
\begin{itemize}
    \item \textbf{Model Scale}: Evaluated on 1.5B parameters; behavior of 7B--72B or reasoning-specialized models (e.g., DeepSeek-R1) remains an open empirical question.
    \item \textbf{Task Scope}: Evaluated on 2-step synthetic arithmetic continuation; does not cover multi-branch logic or interactive theorem proving.
    \item \textbf{Decoding Mode}: Evaluated under unprompted nucleus sampling ($T=0.7, p=0.9$) without external verifier tree search.
    \item \textbf{Causal Boundaries}: The assigned prefix condition is causally identified, but checkpoint comparisons remain observational.
\end{itemize}

\section{Reproducibility \& Artifact Availability}
All code, raw evidence, and verification scripts are archived. The complete neural dataset \texttt{10\_COMPLETE\_QWEN\_PROTOCOL2\_ROLLOUTS.jsonl} (SHA-256: \texttt{b8792f8df375fccdbd800d440cc884ef2db2f860ff3533af01c70161a5cc1d15}) and evaluation package are preserved in the repository.

\section{Conclusion}
We presented \texttt{recovery\_eval}, a state-matched evaluation framework for language-model reasoning trajectories. By eliminating prompt length asymmetry and truncation censoring, we demonstrated that prior claims of a negative post-training recovery interaction were artifacts of evaluation design, and established that 1.5B language models exhibit continuation inertia rather than intrinsic self-correction.

\bibliographystyle{IEEEtran}
\bibliography{references}

\end{document}
